% ECONOMICS_PROBLEM: {"id": "I11-486", "title": "Vertical integration in health care", "jel_code": "I11", "source_id": "OP-0441"}
% !TeX program = xelatex
\documentclass[11pt,a4paper]{article}
\usepackage[margin=23mm]{geometry}
\usepackage{fontspec}
\setmainfont{Latin Modern Roman}
\usepackage{amsmath,amssymb,mathtools}
\usepackage{xcolor,enumitem,booktabs,longtable,array,xurl,hyperref}
\definecolor{muted}{HTML}{53636C}
\hypersetup{colorlinks=true,urlcolor=blue,linkcolor=blue}
\setlength{\parindent}{0pt}
\setlength{\parskip}{5pt}
\newcommand{\R}{\mathbb R}
\newcommand{\E}{\mathbb E}
\newcommand{\N}{\mathbb N}
\newcommand{\ind}{\mathbf 1}
\newcommand{\Var}{\operatorname{Var}}
\newcommand{\Cov}{\operatorname{Cov}}
\newcommand{\supp}{\operatorname{supp}}
\DeclareMathOperator*{\argmin}{arg\,min}
\DeclareMathOperator*{\argmax}{arg\,max}
\newcommand{\entrymeta}[3]{{\sffamily\small\color{muted}\textbf{\texttt{#1}}\quad|\quad #2\par #3\par}}
\newcommand{\Problemref}[1]{\texttt{#1}}
\newcommand{\JELref}[2]{\texttt{#2} (JEL \texttt{#1})}
\begin{document}
\section*{Vertical integration in health care}
\textbf{Problem ID:} \texttt{I11-486}\quad \textbf{JEL:} \texttt{I11}\par
% BEGIN ECONOMICS STATEMENT
\label{problem:OP-0441}
{\sffamily\small\color{muted}Original subject group: Health economics\par}
\label{definitions:problem:30007723}
\textbf{Audit classification:} Published research agenda. Checked source identity and appropriate agenda/background support; expanded intervention, units, population, definedness and causal restrictions. Exact targets and variants are editorial.
\subsection*{Definitions and precise research target}
\textit{Editorial formalization.} Consider Medicare beneficiaries receiving primary care in United States practices acquired by hospitals during 2010–2025. Define \(a\in\{0,1\}\) as hospital ownership versus continued independent ownership of the same practice. For baseline beneficiary \(i\), let \(H_i(a)\) be five-year quality-adjusted survival, \(C_i(a)\) total medical resource cost, and \(T_i(a)\) travel and waiting burden. Fix a value \(\lambda>0\) per quality-adjusted life-year and a monetary valuation \(\rho>0\) per burden unit. The target is
\[\tau_W=E[\lambda\{H_i(1)-H_i(0)\}-\{C_i(1)-C_i(0)\}-\rho\{T_i(1)-T_i(0)\}].\]
Report contrasts by baseline rural residence and clinical risk without conditioning on post-acquisition referral choices. Identify acquisition effects under an explicit ownership counterfactual, retaining patients who move or change providers. Separate coordination changes from billing reclassification and referral steering where additional variation permits; ownership alone does not identify those mechanisms. Include effects on nonacquired practices if spillovers matter and show sensitivity to welfare valuations.

Define quality-adjusted survival as \(H_i(a)=\sum_{t=1}^{5}\beta^t q_{it}(a)\mathbf1\{i\text{ alive at }t\}\), where \(q_{it}\in[0,1]\) is an independently defined health index; deceased persons have zero subsequent survival contribution. \(C_i(a)\) includes all real care costs, including providers changed after acquisition, and \(T_i(a)\) records hours of travel/waiting on the same five-year discount convention. \(\lambda\) has units dollars per discounted quality-adjusted life-year and \(\rho\) dollars per discounted hour. Declare whose time and costs are included and require finite means. Site-of-service reimbursement changes alone are not real cost changes. Ownership assignment is at practice/system level; a treated-practice randomization or justified counterfactual must address endogenous acquisition and spillovers. The expression is a declared net-health-benefit objective, not an identified utilitarian welfare sum without further preferences. All expectations use a stated baseline population and probability law. Identification means every admissible data-generating model with the same observed-data law yields the same displayed target; otherwise report the identified set. Exact equations, horizons and assumptions are editorial, rather than source-given canonical conjectures.
\subsection*{Variants and assumption changes}
\begin{enumerate}
\item \textbf{Editorial saturation variant.} Include nonacquired practices and patients in market welfare \(W(s)\), where \(s\) is ownership prevalence. Compare with isolated-practice contrasts.
\item \textbf{Editorial payment-neutral variant.} Cross acquisition with removal of site-of-service reimbursement advantages. The interaction identifies this channel only under credible variation in both regimes.
\end{enumerate}
\subsection*{References and source audit}
\textbf{1.} Consolidation quality and access evidence gaps verified; claimed exact geographic-proximity passage not independently retrieved. Locator: RAND RR-A1820-1, 30 September 2022; report-page Key Takeaways. Report metadata and primary RAND summary inspected; PDF follow-up retrieval timed out. The supplied p.13 proximity locator remains unverified. URL: \url{https://www.rand.org/content/dam/rand/pubs/research_reports/RRA1800/RRA1820-1/RAND_RRA1820-1.pdf}.
\begin{enumerate}\raggedright
\item\hypertarget{definitions:source:30007723:1}{} Jodi L. Liu; Zachary M. Levinson; Annetta Zhou; Xiaoxi Zhao; PhuongGiang Nguyen; Nabeel Qureshi. 2022. \emph{Environmental Scan on Consolidation Trends and Impacts in Health Care Markets}. RAND RR-A1820-1, 30 September 2022; report-page Key Takeaways.
\url{https://www.rand.org/content/dam/rand/pubs/research_reports/RRA1800/RRA1820-1/RAND_RRA1820-1.pdf}.
DOI: \url{https://doi.org/10.7249/RRA1820-1}.
\end{enumerate}

\subsection*{Common conventions: Source mathematical conventions}

These conventions apply to each entry unless explicitly replaced.

\textbf{Models and identification.} A primitive model \(m\) specifies all probability laws, feasible choices, information, equilibrium and measurement rules used by its target. The admissible class \(\mathcal M\) is fixed independently of outcomes. If \(P_O(m)\) is its observable probability law and \(\tau(m)\) the target, its identified set at observable law \(P\) is
\[
 \mathcal I_\tau(P)=\{\tau(m):m\in\mathcal M,\ P_O(m)=P\}.
\]
Point identification means that this set is a singleton. An empty set signals incompatibility of data and assumptions. Coordinate bounds need not describe the joint set. Bounds are sharp only if every retained target is attainable or arbitrarily approachable by compatible models. Finite bounds require explicit restrictions on unobserved quantities; unrestricted confounding, hidden wealth, extrapolation or arbitrary equilibrium selection can yield uninformative sets.

\textbf{Causal interventions.} A potential outcome \(Y_i(p)\) is the outcome under a complete policy regime \(p\), including timing, financing, information and market spillovers. The target population and units are fixed before treatment. With interference, use the complete assignment vector or a justified exposure mapping; individual no-interference notation alone cannot identify an economy-wide intervention. A difference of mean potential outcomes does not require the cross-world joint distribution. Conditioning on post-treatment migration, survival, enrollment or employment changes the estimand and needs separate justification.

\textbf{Statistical statements.} An estimator, confidence set, forecast or test specifies a sampling law, model class and loss/coverage criterion. Uniform \((1-\alpha)\)-coverage means \(\inf_{m\in\mathcal M}\Pr_m\{\tau(m)\in C_n\}\geq1-\alpha\), or an explicitly asymptotic version. The whole parameter space trivially has coverage, so informativeness requires a stated width, risk or power objective. A forecasting improvement is relative to a named benchmark, information set and evaluation population.

\textbf{Equilibrium and optimization.} An equilibrium comprises feasible plans, optimality, beliefs where needed, budget constraints and all market/resource clearing. The primitives or a stated benchmark define each correspondence \(\mathcal E(p;m)\). Nonemptiness is assumed or proved, never inferred just from compact policy sets. A selected-equilibrium target uses a declared selection rule; a robust target quantifies over all permitted equilibria. Use suprema or infima unless attainment is established. An \(\varepsilon\)-optimal policy is feasible and within \(\varepsilon\) of the finite optimum value. Report an empty feasible set separately.

\textbf{Welfare and accounting.} Normative utility, interpersonal normalization, social weights, numeraire, discounting and the use of public receipts are primitives. Behavioral utility need not coincide with the welfare criterion. Household welfare includes ownership income and rents once; adding profits again to compensating variation generally double-counts them. A monetary equivalent is defined by an explicit transfer and price convention, with monotonicity and range conditions. Average effects, mechanism contributions, transfers and welfare losses are different targets. Infinite sums require convergence and dynamic feasible plans require terminal or no-Ponzi conditions.

\textbf{Source versions.} A working-paper date, revision date and journal publication year are distinguished. A DOI for a journal version is not automatically the DOI of an earlier manuscript. Locators refer to the stated version and printed pages where specified. A metadata-only check does not verify a claim about a full-text conclusion. Every entry's source subsection narrows the provenance claim accordingly.
% END ECONOMICS STATEMENT
\end{document}
